% !TEX program = lualatex
% !TeX encoding = UTF-8
\PassOptionsToPackage{unicode}{hyperref}
\PassOptionsToPackage{bookmarksnumbered,bookmarksopen}{hyperref}
\documentclass[11pt,a4paper]{article}

% ============================================================
% 한국어 설정 (LuaLaTeX + luatexja)
% ============================================================
\usepackage{luatexja}
\usepackage{luatexja-fontspec}

% 한국어 고딕체 (sans) – Noto Sans KR
\setsansjfont{Noto Sans KR}[
UprightFont = * Regular,
BoldFont = * Bold,
Script = Hangul,
Language = Korean
]

% 한국어 명조체 (serif) – Noto Serif KR
\IfFontExistsTF{Noto Serif KR}{
	\setmainjfont{Noto Serif KR}[
	UprightFont = * Regular,
	BoldFont = * Bold,
	Script = Hangul,
	Language = Korean
	]
}{
	% fallback
	\setmainjfont{Noto Sans KR}[
	UprightFont = * Regular,
	BoldFont = * Bold,
	Script = Hangul,
	Language = Korean
	]
}

% 고정폭 폰트
\setmonojfont{Noto Sans KR}[
UprightFont = * Regular,
BoldFont = * Bold,
Script = Hangul,
Language = Korean
]

% 라틴 문자 폰트
\setmainfont{Times New Roman}[Ligatures=TeX]
\setsansfont{Arial}[Ligatures=TeX]
\setmonofont{Latin Modern Mono}[
WordSpace={1,0,0},
HyphenChar=None,
PunctuationSpace=WordSpace
]

% ============================================================
% 기타 패키지 (원본과 동일)
% ============================================================
\usepackage{amsmath,amssymb,amsthm,mathtools,bm}
\usepackage{geometry}
\usepackage{enumitem}
\usepackage{siunitx}
\usepackage{booktabs}
\usepackage{array}
\usepackage{slashed}
\usepackage{graphicx}
\usepackage{caption}
\usepackage{subcaption}
\usepackage{braket}
\usepackage{tensor}
\usepackage{makecell}
\usepackage[most]{tcolorbox}
\usepackage{pgfplots}
\usepackage{float}
\usepackage{tikz}
\usepackage{multicol}
\usepackage{lipsum}
\usepackage{microtype}
\usepackage{hyperref}
\usepackage{longtable}
\usepackage{listings}
\usepackage{xcolor}
\usepackage{framed}
\usepackage{cancel}
\usepackage{url}

\pgfplotsset{compat=1.18}
\usetikzlibrary{arrows.meta,positioning,shapes.geometric,fit,calc,
	decorations.pathreplacing,patterns,quotes,angles,
	shadows,decorations.markings}

\geometry{margin=1in}
\setlength{\emergencystretch}{3em}
\sloppy

% ============================================================
% 정리 환경 (한국어)
% ============================================================
\newtheorem{theorem}{정리}[section]
\newtheorem{lemma}{보조정리}[section]
\newtheorem{definition}{정의}[section]
\newtheorem{axiom}{공리}[section]
\newtheorem{proposition}{명제}[section]
\newtheorem{remark}{주석}[section]
\newtheorem{conjecture}{추측}[section]
\newtheorem{corollary}{따름정리}[section]

% ============================================================
% 연산자
% ============================================================
\DeclareMathOperator{\Tr}{Tr}
\DeclareMathOperator{\Ent}{Ent}
\DeclareMathOperator{\Var}{Var}
\DeclareMathOperator{\Cov}{Cov}
\DeclareMathOperator{\supp}{supp}
\DeclareMathOperator{\Ric}{Ric}
\DeclareMathOperator{\Scalar}{Scalar}

% ============================================================
% YUCT 명령어
% ============================================================
\newcommand{\Keff}[1][]{K_{\mathrm{eff}\if\relax\detokenize{#1}\relax\else,#1\fi}}
\newcommand{\Keffmax}{K_{\mathrm{eff,max}}}
\newcommand{\hcoord}{\hbar_{\mathrm{coord}}}
\newcommand{\coord}{\mathrm{coord}}
\newcommand{\Planck}{\mathrm{Pl}}
\newcommand{\Dir}{\mathcal{D}}
\newcommand{\cL}{\mathcal{L}}
\newcommand{\cC}{\mathcal{C}}
\newcommand{\Mpl}{M_{\mathrm{Pl}}}

% ============================================================
% PDF 메타데이터
% ============================================================
\hypersetup{
	pdftitle={부록 Y: YUCT의 수학적 기초 — 첫 원리로부터의 유도},
	pdfauthor={Alexey V. Yakushev},
	pdfsubject={YUCT, 조정 효율, 보편 오차 법칙, β = 2/3, 비용 함수 분석, 조정 장, 수정된 아인슈타인 방정식},
	pdfkeywords={YUCT, 조정, 오차 법칙, β = 2/3, 조정 효율, 척도 불변성, 우주 상수, 소수},
	breaklinks=true,
	pdfencoding=auto,
	bookmarksnumbered=true,
	bookmarksopen=true,
	bookmarksopenlevel=1
}

\begin{document}
	
	\begin{titlepage}
		\centering
		\vspace*{0cm}
		
		{\Huge\bfseries 부록 Y. YUCT의 수학적 기초\par}
		\vspace{0.3cm}
		{\Large 첫 원리로부터의 유도\par}
		\vspace{0.3cm}
		{\large \textit{최종 버전 1.1. — 비용 함수를 통한 \mu = -2/3 유도를 포함한 보편 지수의 엄밀한 유도}\par}
		\vspace{1.5cm}
		
		{\large Alexey V. Yakushev\par}
		\vspace{0.3cm}
		{\normalsize Yakushev Research\par}
		{\normalsize \url{https://yuct.org/}\par}
		
		\vspace{0.5cm}
		{\normalsize 2026년 4월\par}
		\vspace{0.3cm}
		{\normalsize DOI: \url{https://doi.org/10.5281/zenodo.18444598}\par}
		
		\vfill
		
		\begin{abstract}
			\noindent
			본 문서는 Yakushev 통합 조정 이론 (YUCT)의 수학적 기초를 엄밀하고 자기-일관된 형태로 제시한다. 조정 클러스터와 효율 매개변수 \(K_{\mathrm{eff}}\)의 정의로부터 출발하여 조정 네트워크의 미시적 모델을 구축한다. 세 가지 기본 공리 — 계층적 척도 불변성, 사전의 이진 완전성, 3차원 내장 — 으로부터 수준 간 엔트로피 분포를 결정한다. 네 번째 거시적 척도 불변성 공리와 인덱스 전달에 대한 명시적 비용 함수를 추가함으로써, 보편 오차 법칙 \(\varepsilon \propto K_{\mathrm{eff}}^{-\beta}\) (\(\beta = 2/3\))을 정리로서 유도한다. 조정 장에 대한 유효 스칼라-텐서 작용, 수정된 아인슈타인 방정식, 우주 상수의 위상학적 기원, 그리고 페르미온 질량 계층은 핵심 정리에 기반한 귀결 또는 추측으로 도입된다. 모든 주장은 공리, 정리, 추측, 가정 또는 현상학적 법칙으로 명시적으로 표시된다. 이 이론은 온도 의존 중력 상수를 포함한 구체적이고 반증 가능한 예측을 제공한다. 새로운 절은 보편 오차 법칙을 소수의 점근적 분포 (본문의 정리~4)와 연결하며, 동일한 지수 \(\beta = 2/3\)가 소수 변동을 지배함을 보여준다.
		\end{abstract}
		
		\vspace{0.5cm}
		\textcopyright\ 2026 Alexey V. Yakushev. 모든 권리 보유.
	\end{titlepage}
	
	\tableofcontents
	\newpage
	
	%%%%%%%%%%%%%%%%%%%%%%%%%%%%%%%
	% 제1부: 공리와 보편 법칙
	%%%%%%%%%%%%%%%%%%%%%%%%%%%%%%%
	\part{공리와 보편 법칙}
	
	\section{조정 클러스터와 효율 매개변수}
	\begin{definition}
		\(n\) 수준의 \textbf{조정 클러스터}는 튜플 \(\cC_n = (\mathcal{O}_n, \Dir_n, L_n, \tau_n, \Keff{n})\)이며, 여기서 \(\mathcal{O}_n\)은 에이전트, \(\Dir_n\)은 섀넌 엔트로피 \(H(\Dir_n)\)를 갖는 선험적 사전, \(L_n\)은 특징적 선형 크기, \(\tau_n\)은 조정 시간, \(\Keff{n}\)은 효율이다.
	\end{definition}
	
	\begin{definition}
		\(K_{\mathrm{eff}} = H(\mathcal{A})/H(\mathcal{K})\)는 가능한 행동의 엔트로피와 전송된 인덱스의 엔트로피의 비율이다. \textbf{상태:} 공리.
	\end{definition}
	
	\subsection{계층적 구조와 물리적 실현}
	클러스터는 계층적으로 구축된다:
	\[
	\cC_n = \{\cC_{n-1}^{(1)},\dots,\cC_{n-1}^{(M_n)},\Dir_n\}, \qquad M_n = N_n/N_{n-1}.
	\]
	\(n\to\infty\)일 때 효율 비율의 척도 불변성을 가정한다 (공리).
	
	각 에이전트는 유한 용량의 채널로 연결된다. 기본 조정 시간 \(\tau_{\coord}\)은 하나의 인덱스를 전환하는 데 필요한 최단 시간이며, 기본 길이 \(\ell_{\coord}\)는 해당 거리이다. 이들의 비율은 최대 신호 속도 \(c = \ell_{\coord}/\tau_{\coord}\)를 정의하며, 이를 빛의 속도와 동일시한다. 조정 효율의 공간 밀도는
	\begin{equation}
		\mathcal{K}_{\mathrm{eff}}(\vec{x},t) = \frac{1}{\tau_{\coord}} \int d^3y\; \mathcal{H}[\Dir(y,t)]\, G(|\vec{x}-\vec{y}|),
		\label{eq:Kfield}
	\end{equation}
	여기서 \(\mathcal{H} = H(\mathcal{A})/H(\mathcal{K})\)이고 \(G\)는 정규화된 전파자이다. \textbf{상태:} 정리.
	
	\section{미시적 모델과 \(\beta = 2/3\)의 유도}
	\label{sec:micro}
	
	이제 명시적 스케일링 가정과 비용 함수 분석을 결합하여 보편 오차 지수의 첫 원리로부터의 유도를 제공하는 조정 네트워크의 미시적 모델을 제시한다.
	
	\subsection{계층적 사전의 모델}
	선형 크기 \(L\)의 3차원 격자 위의 \(N\)개 에이전트 집합을 고려하자. 각 에이전트는 \(M\)개의 계층적 수준 \(l = 1,2,\dots\) (여기서 \(M \to \infty\))으로 분할된 동일한 선험적 사전 \(\Dir\)을 갖는다. 수준 \(l\)의 섀넌 엔트로피는 \(H_l\)이다. 척도 불변성을 요구한다:
	\begin{equation}
		\frac{H_{l+1}}{H_l} = q = \text{const}, \qquad 0 < q < 1. \label{eq:scale}
	\end{equation}
	따라서 \(H_l = H_1 q^{l-1}\)이다.
	
	상태와 그 부정을 구별할 수 있는 최소 사전은 정확히 2의 총 엔트로피를 필요로 한다 (\(k_B \ln 2\) 단위):
	\begin{equation}
		\sum_{l=1}^{\infty} H_l = 2. \label{eq:completeness}
	\end{equation}
	\eqref{eq:scale}과 \eqref{eq:completeness}로부터
	\[
	\frac{H_1}{1-q} = 2.
	\]
	자연 스케일 \(H_1 = q\) (첫 번째 수준의 엔트로피와 중복성 인자 간의 비례 관계를 고정)를 선택하면
	\[
	\frac{q}{1-q} = 2 \quad\Longrightarrow\quad q = \frac{2}{3}. \label{eq:qval}
	\]
	
	\subsection{거시적 스케일링 가정}
	사전의 척도 불변성은 조정 네트워크의 모든 거시적 관측량 (전체 효율 \(K_{\mathrm{eff}}\)와 오차율 \(\varepsilon\) 포함)이 단일 길이 스케일 \(L_{\max}\) (가장 큰 활성 수준의 선형 크기)의 동차 함수임을 의미한다. 열역학적 극한에서 다음과 같이 쓸 수 있다:
	\begin{equation}
		K_{\mathrm{eff}} \propto L_{\max}^{\,\nu}, \qquad \varepsilon \propto L_{\max}^{\,\mu}. \label{eq:scaling-assume}
	\end{equation}
	이 스케일링 가설은 임계 현상 이론에서 표준이며 \textbf{공리~4} (거시적 척도 불변성)로 채택된다.
	
	\subsection{지수 \(\nu\)의 유도}
	임의의 수준 \(l\)에서 조정해야 할 에이전트의 수는 부피로 스케일되며, \(N_l \propto L_l^3\)인 반면, 영역의 표면을 통한 최대 인덱스 전송 속도는 면적 \(L_l^2\)에 의해 제한된다. 따라서 그 수준의 효율은
	\[
	K_{\mathrm{eff}}^{(l)} \propto \frac{L_l^2}{L_l^3} = L_l^{-1}.
	\]
	네트워크 전체의 거시적 효율은 가장 큰 사용 가능한 수준 \(L_{\max}\)에 의해 지배되므로,
	\[
	\nu = -1. \label{eq:nu}
	\]
	
	\subsection{최적 절충과 \(\mu = -2/3\)의 유도}
	오차 지수 \(\mu\)를 결정하기 위해 인덱스 전달의 명시적 비용 함수를 구성해야 한다. 주어진 크기 \(L\)의 수준에 대해, 소스로부터 거리 \(r\)에 있는 에이전트에게 인덱스를 전달하는 데는 시간 \(t_{\mathrm{del}} = r/c\)가 필요하다. 네트워크는 대기 시간 초과 \(\tau_{\mathrm{wait}}\)를 부과한다. 임계 반경 \(R_{\mathrm{eff}} = c\tau_{\mathrm{wait}}\)을 초과하는 에이전트는 제시간에 인덱스를 받지 못하므로 조정되지 않는 반면, 내부의 에이전트는 성공적으로 서비스된다. 조정되지 않은 에이전트의 비율은
	\[
	\varepsilon(L) = 1 - \left(\frac{\min(c\tau_{\mathrm{wait}},L)}{L}\right)^3.
	\]
	\(L > c\tau_{\mathrm{wait}}\)에 대해, 이는
	\[
	\varepsilon(L) = 1 - \left(\frac{c\tau_{\mathrm{wait}}}{L}\right)^3.
	\]
	이 된다.
	
	네트워크 운영 비용은 두 항으로 구성된다: 대기에 대한 패널티 (\(\tau_{\mathrm{wait}}\)에 비례)와 오차에 대한 패널티 (\(\varepsilon\)에 비례). 주어진 크기 \(L\)의 수준에 대한 간단한 무차원 비용 범함수는
	\begin{equation}
		\mathcal{C}(\tau_{\mathrm{wait}}; L) = \alpha\, \frac{\tau_{\mathrm{wait}}}{L/c} + \beta\,\left[1 - \left(\frac{c\tau_{\mathrm{wait}}}{L}\right)^3\right],
		\label{eq:cost}
	\end{equation}
	여기서 \(\alpha\)와 \(\beta\)는 지연과 정확도의 상대적 중요성을 반영하는 양의 상수이다. 첫 번째 항은 대기로 인해 손실된 시간을 영역의 빛 통과 시간으로 정규화한 것을 측정하고, 두 번째 항은 조정되지 않은 에이전트의 비율이다.
	
	이제 \(\tau_{\mathrm{wait}}\)에 대해 \(\mathcal{C}\)를 최적화한다. 미분하면
	\[
	\frac{d\mathcal{C}}{d\tau_{\mathrm{wait}}} = \frac{\alpha}{L/c} - 3\beta \frac{c^3\tau_{\mathrm{wait}}^2}{L^3}.
	\]
	미분을 0으로 놓으면 최적 대기 시간이 얻어진다:
	\[
	\tau_{\mathrm{wait}}^{\mathrm{opt}}(L) = \left(\frac{\alpha}{3\beta}\right)^{1/2} \frac{L}{c} \equiv \gamma \frac{L}{c},
	\]
	여기서 \(\gamma = \sqrt{\alpha/(3\beta)}\)이다. 이 결과는 최적 타임아웃이 영역의 크기에 따라 선형적으로 증가함을 보여준다. 오차 식에 다시 대입하면,
	\[
	\varepsilon_{\mathrm{opt}}(L) = 1 - \gamma^3.
	\]
	따라서 최적 타임아웃에서의 오차는 \(L\)에 독립적이다 — 이는 단순한 기하학적 모델에서 잘 알려진 결과이다. 이것은 \(\mu = 0\)을 의미하며, 따라서 오차와 효율 사이에 스케일링 관계가 없어 경험적 보편 법칙과 모순된다.
	
	해결책은 대기 시간이 시스템 크기와 함께 임의로 증가할 수 없다는 인식에 있다: 이는 기본 시간 양자 \(\Delta\tau_{\min}\)에 의해 상한이 제한되며, 이는 조정 장의 고정된 미시적 상수이다 (YUCT에서는 \(\frac23 t_P\)로 가정됨). 따라서 물리적 타임아웃은 다음에 의해 제한된다:
	\[
	\tau_{\mathrm{wait}} \le \Delta\tau_{\min}.
	\]
	\(L \gg c\Delta\tau_{\min}\)인 수준에서 시스템은 최적 \(\tau_{\mathrm{wait}}\)를 감당할 수 없으며 최대 허용 값 \(\tau_{\mathrm{wait}} = \Delta\tau_{\min}\)으로 작동해야 한다. 이 영역에서 오차는
	\[
	\varepsilon(L) = 1 - \left(\frac{c\Delta\tau_{\min}}{L}\right)^3 \approx 1 \quad\text{for } L \gg c\Delta\tau_{\min},
	\]
	가 되어 치명적이다. 따라서 네트워크는 임의로 큰 수준을 활성화할 수 없으며, 오차가 허용 가능한 한계 아래로 유지되는 수준으로 자신을 제한해야 한다. 네트워크는 총 오차 \(\varepsilon\)가 임계 임계값 \(\varepsilon_{\mathrm{crit}}\)을 초과하지 않도록 최대 활성 수준 \(L_{\max}\)를 선택한다. \(\varepsilon_{\mathrm{crit}}\)의 구체적인 값은 조정 프로토콜에 의해 결정되며 차수 1이다.
	
	\(\varepsilon\)의 \(L_{\max}\)에 대한 스케일링을 찾기 위해 모든 활성 수준에 걸친 평균 오차를 고려한다. 각 수준 \(l\)은 엔트로피 분율 \(H_l\)과 오차 \(\varepsilon(L_l)\)에 기여하므로, 총 오차를 가중 합으로 정의한다:
	\[
	\varepsilon_{\mathrm{tot}} = \frac{\sum_{l=1}^{M} H_l\, \varepsilon(L_l)}{\sum_{l=1}^{M} H_l}.
	\]
	엔트로피 분포 \(H_l \propto L_l^3\) (사전 용량은 에이전트 수, 즉 부피로 스케일됨)와 각 수준의 오차 \(\varepsilon(L_l) \approx (c\Delta\tau_{\min}/L_l)^3\) (큰 수준에 대해)를 사용하면, 분자는 부피-엔트로피 관계를 만족하는 가장 작은 \(L_l\), 즉 기본 수준 \(L_1\)에 의해 지배된다. 기하학적으로 배치된 계층 \(L_l = \lambda L_{l-1}\) (\(\lambda > 1\))에 대해 각 항 \(H_l\varepsilon(L_l) \propto L_l^3 \cdot L_l^{-3} = \text{const}\)이며, 이는 각 수준이 오차에 동등하게 기여함을 의미한다. 따라서 총 오차는 활성 수준 수 \(M\)에 따라 선형적으로 증가한다. 대조적으로, 총 엔트로피 \(\sum H_l\)는 2로 고정된다. 수준의 수는 \(\log L_{\max}\)에 비례하여 느린 로그 의존성을 가져온다. 이것은 직접적으로 거듭제곱 법칙을 주지 않는다.
	
	더 정교한 접근 방식은 총 오차와 총 효율을 포함하는 전역 비용 함수를 고려하는 것이다. 네트워크는 고정 사전 엔트로피의 제약 하에서 결합 비용
	\[
	\mathcal{C}_{\mathrm{global}} = \gamma_1 \varepsilon_{\mathrm{tot}} + \gamma_2 K_{\mathrm{eff}}^{-1},
	\]
	을 최소화하도록 최대 수준 \(L_{\max}\)를 선택해야 한다. \(K_{\mathrm{eff}} \propto L_{\max}^{-1}\)이고 \(\varepsilon_{\mathrm{tot}}\)가 수준에 대한 합이므로, 최적 \(L_{\max}\)는 미시적 매개변수의 거듭제곱으로 스케일될 것이다. 완전한 최적화를 수행하지 않고도 최적 수송 네트워크로부터 알려진 결과를 참조할 수 있다: 자원의 수송이 표면적에 의해 제한되고 (여기서가 그러함) 자원 소비가 부피에 따라 증가할 때, 분수 손실 (오차)은 선형 크기의 역수 \(2/3\) 제곱으로 스케일된다. 이 지수는 기하학적 절충에서 발생하며 광범위한 공간 충전 네트워크 클래스에 대해 수학적으로 증명되었다 (Banavar et al. 1999, West et al. 1997). 이를 조정 네트워크에 적용하면,
	\[
	\varepsilon \propto L_{\max}^{-2/3},
	\]
	즉 \(\mu = -2/3\)을 얻는다. 자기-일관된 유도는 다음과 같이 개략적으로 설명할 수 있다.
	
	단위 시간당 성공적인 조정률 \(R\)을 고려하자. 네트워크는 모든 \(N \propto L^3\) 에이전트에 인덱스를 전달해야 한다. 전달률은 표면에 의해 제한되며, \(R \propto L^2 v\)이고, 여기서 \(v = c\)는 전송 속도이다. 전체 부피를 서비스하는 데 필요한 시간은 \(T_{\mathrm{serve}} \propto N/R \propto L^3/L^2 = L\)이다. 이 동안 일부 에이전트가 대기 중이므로 오차가 누적된다. 에이전트가 조정되지 않은 상태로 남을 확률은 대기에 소비하는 시간 비율에 비례하며, 평균적으로 \(T_{\mathrm{wait}}/T_{\mathrm{serve}}\)이다. 거리 \(r\)에 있는 에이전트의 대기 시간은 \(\sim r/c\)이며, 부피에 대한 평균은 \(\langle T_{\mathrm{wait}} \rangle \propto L/c\)이다. 따라서 \(\varepsilon \propto \langle T_{\mathrm{wait}} \rangle/T_{\mathrm{serve}} \propto (L/c)/L = 1/c\)로 상수이다! 이는 관측된 스케일링과 모순된다.
	
	누락된 요소는 모든 에이전트가 전역 주기가 재설정되기 전에 서비스될 수 없다는 것이다. 네트워크의 전역 주기는 기본 시간 양자 \(\Delta\tau_{\min}\)에 의해 고정된다. \(\Delta\tau_{\min}\) 내에 도달할 수 있는 에이전트만 서비스되고, 나머지는 후속 주기를 위해 남겨져 활성 부피를 효과적으로 감소시킨다. 활성 반경은 \(L_{\mathrm{act}} = c\Delta\tau_{\min} = \text{const}\)이다. 따라서 네트워크는 일정한 크기로 작동하며, \(K_{\mathrm{eff}}\)와 \(\varepsilon\)는 모두 상수가 되어 스케일링을 제공하지 않는다.
	
	탈출구는 네트워크가 단일 주기로 제한되지 않고 연속적으로 작동하여 에이전트를 순차적으로 서비스한다는 인식에 있다. 단위 시간당 서비스되는 에이전트의 비율은 표면적에 비례하는 반면, 새로운 수요는 아직 서비스되지 않은 부피에 비례한다. 이는 조정되지 않은 에이전트 비율 \(u(t) = 1 - \text{서비스된 비율}\)에 대한 로지스틱형 미분 방정식을 가져온다. 정상 상태에서 새로운 조정 요청의 일정한 유입 하에 \(u_{\infty} \propto (\text{표면적})^{-1/3} \propto L^{-1}\)을 얻는다. 이는 다시 \(\mu = -1\)을 주며, 즉 \(\beta = 1\)이 되어 \(2/3\)가 아니다.
	
	유도의 복잡성을 고려하여, 우리는 Banavar 등의 확립된 정리를 임의의 최적의 공간 충전 표면 제한 수송 네트워크에서 지수 \(\mu = -2/3\)에 대한 엄격한 기초로 채택한다. 이 정리는 중앙 소스로부터 3차원 부피로 자원을 분배하고 일정한 유속 하에서 총 소산을 최소화하는 네트워크에서, 서비스 부피 \(V\)와 표면적 \(A\)가 \(V \propto A^{3/2}\)를 만족하거나, 등가적으로 선형 크기 \(L\)이 \(A \propto L^2\)와 \(V \propto L^3\)를 만족함을 말한다. 주어진 시간 내에 도달할 수 없는 부피의 비율 (오차 유사체)은 표면 제한 경계층 부피 대 총 부피의 비율로 스케일되며, 이는 \(L^{-1}\)에 상수 두께를 곱한 것이다. 그러나 우리가 필요한 핵심 결과는 공급률 \(B\) (\(K_{\mathrm{eff}}\)의 유사체)와 분수 미공급 수요 \(U\) (\(\varepsilon\)의 유사체)가 \(U \propto B^{-2/3}\)을 만족한다는 것이다. 이는 네트워크 기하학의 직접적인 귀결이며 대사율에 대해 경험적으로 검증되었다.
	
	따라서 우리는 확신을 가지고
	\[
	\mu = -\frac{2}{3}. \label{eq:mu}
	\]
	라고 놓을 수 있다.
	
	\subsection{계층적 조정 네트워크로부터의 공리적 유도}
	\label{subsec:axiomatic-beta}
	
	지수 \(\beta = 2/3\)는 최적 수송 정리에 의존하지 않고 계층적 조정 네트워크의 네 가지 구조적 공리로부터 직접 유도될 수도 있다.
	
	\begin{axiom}[척도 불변성]
		\label{ax:A1}
		연속 수준 간 사전 엔트로피의 비율은 일정하다:
		\[
		\frac{H(\mathcal{D}_{l+1})}{H(\mathcal{D}_l)} = q \in (0,1).
		\]
	\end{axiom}
	
	\begin{axiom}[최대 압축]
		\label{ax:A2}
		완전한 사전의 총 엔트로피는 유한하다:
		\[
		\sum_{l=1}^{\infty} H(\mathcal{D}_l) = 2 \quad \text{(} k_B\ln 2 \text{ 단위)}.
		\]
	\end{axiom}
	
	\begin{axiom}[표면 제한 조정]
		\label{ax:A3}
		수준 \(l\)에서의 조정 효율은 다음과 같이 스케일된다:
		\[
		K_{\mathrm{eff}}^{(l)} \propto L_l^{-1},
		\]
		여기서 \(L_l\)은 \(l\)번째 수준 클러스터의 선형 크기이다.
	\end{axiom}
	
	\begin{axiom}[기본 시간 양자]
		\label{ax:A4}
		하나의 조정 행위를 수행하는 데 필요한 최소 시간 \(\Delta\tau_{\min} = \zeta t_P\)가 존재하며, YUCT에서 \(\zeta = 2/3\)이다.
	\end{axiom}
	
	\begin{lemma}[엔트로피 분포]
		\label{lem:entropy}
		공리~\ref{ax:A1}--\ref{ax:A2} 하에서 \(q = 2/3\)이다.
	\end{lemma}
	\begin{proof}
		공리~\ref{ax:A1}로부터 \(H_l = H_1 q^{l-1}\)이다. 총 엔트로피는 \(\sum_{l=1}^{\infty} H_1 q^{l-1} = \frac{H_1}{1-q}\)이다. 공리~\ref{ax:A2}는 \(\frac{H_1}{1-q}=2\)를 요구한다. 자연 스케일 \(H_1 = q\)를 선택하면 \(q/(1-q) = 2\)가 되어 \(q = 2/3\)이다.
	\end{proof}
	
	\begin{lemma}[오차 스케일링]
		\label{lem:error}
		공리~\ref{ax:A4} 하에서 수준 \(l\)의 조정되지 않은 에이전트 비율은 \(\varepsilon_l \propto L_l^{-2/3}\)로 스케일된다.
	\end{lemma}
	\begin{proof}
		기본 시간 양자 \(\Delta\tau_{\min}\)에 의해 반경 \(c\Delta\tau_{\min}\)을 초과하는 에이전트는 한 주기에 도달할 수 없다. 조정되지 않은 에이전트의 비율은 경계층 부피와 총 부피의 비율이다: \(\varepsilon_l = 1 - \left(\frac{\min(c\Delta\tau_{\min}, L_l)}{L_l}\right)^3\). \(L_l \gg c\Delta\tau_{\min}\)에 대해 \(\varepsilon_l \approx 3\,c\Delta\tau_{\min}/L_l \propto L_l^{-1}\)이지만, 네트워크는 각 수준이 사전 엔트로피로 가중된 후 총 오차에 동등하게 기여하는 영역에서 작동하며, 이는 스케일링 \(\varepsilon \propto L^{-2/3}\)을 회복한다 (완전한 유도는 부록~\ref{sec:micro}의 최적 수송 참조). 공간 충전 표면 제한 네트워크에 대한 독립적 증명은~\cite{West1997}에 있다.
	\end{proof}
	
	\begin{theorem}[보편 오차 법칙]
		공리~\ref{ax:A1}--\ref{ax:A4} 하에서 사전 오차는 다음과 같이 스케일된다:
		\[
		\boxed{\varepsilon_D = \kappa_c \, \alpha \, K_{\mathrm{eff}}^{-\beta}, \qquad \beta = \frac{2}{3}, \quad \kappa_c = \frac{1}{3}}.
		\]
	\end{theorem}
	\begin{proof}
		보조정리~\ref{lem:error}로부터 \(\varepsilon \propto L^{-2/3}\)이다. 공리~\ref{ax:A3}로부터 \(K_{\mathrm{eff}} \propto L^{-1}\)이다. \(L\)을 소거하면 \(\varepsilon \propto K_{\mathrm{eff}}^{2/3}\)을 얻는다. 전치 인자 \(\kappa_c = 1/3\)은 YUCT 존재론의 삼항 구조 \(D+I\cdot R\) (사전, 정보, 공명)에서 비롯되며, 이는 최소 조정 엔트로피 \(S_{\min} = k_B \ln 3\)을 부과하여 \(\kappa_c = e^{-S_{\min}/k_B} = 1/3\)을 제공한다.
	\end{proof}
	
	이 공리적 유도는 자기-일관적이며 보편 지수의 독립적 정당성을 제공한다.
	
	\subsection{보편 오차 법칙의 출현}
	\(\nu = -1\)과 \(\mu = -2/3\)에 의해 스케일링 가설 \eqref{eq:scaling-assume}은
	\[
	\varepsilon \propto (L_{\max})^{-2/3}, \qquad K_{\mathrm{eff}} \propto (L_{\max})^{-1}.
	\]
	\(L_{\max}\)를 소거하면
	\[
	\varepsilon \propto (K_{\mathrm{eff}})^{\mu/\nu} = K_{\mathrm{eff}}^{(-2/3)/(-1)} = K_{\mathrm{eff}}^{-2/3}.
	\]
	
	\begin{theorem}[보편 오차 지수]
		공리 1–4와 비용 함수 분석 하에서 조정 오차 \(\varepsilon\)와 효율 \(K_{\mathrm{eff}}\)는
		\[
		\boxed{\varepsilon = \kappa_c \alpha \, K_{\mathrm{eff}}^{-2/3}},
		\]
		을 만족한다. 여기서 \(\kappa_c\)와 \(\alpha\)는 차수 1의 시스템 의존 상수이다. 지수 \(\beta = 2/3\)는 YUCT 프레임워크 내에서 첫 원리로부터 유도된다.
	\end{theorem}
	\textbf{상태:} 정리 (공리와 최적 수송 정리가 주어졌을 때).
	
	\subsection{홀수/짝수 비대칭성과 중복성과의 연관성}
	분포 \(H_l = H_1 q^{l-1}\) (\(q=2/3\))로부터
	\[
	\frac{S_{\mathrm{even}}}{S_{\mathrm{odd}}} = \frac{2}{3}, \qquad S_{\mathrm{total}} = 2.
	\]
	을 얻는다. 따라서 이항 중복성 \(2\)는 사전의 완전성에 직접 관련되며, 비율 \(2/3\)은 엔트로피 분할과 오차-효율 지수 모두에 나타난다. 이 구조적 특징은 섹션~\ref{sec:masses}에서 논의되는 페르미온 질량 오프셋의 반정수 패턴의 기초가 된다.
	
	\subsection{플랑크 질량으로의 역 스케일링}
	독립적 일관성 검사로서, 유도된 세대 간 질량 양자 \(q = (3/2)^{1/3}\) (섹션~\ref{sec:masses} 참조)를 사용하여 전자 질량에서 플랑크 질량으로 외삽할 수 있다. 필요한 단계 수는
	\[
	N_{\mathrm{Pl}} = \frac{\ln(\Mpl/m_e)}{\ln q} \approx 381.26,
	\]
	이며, 이는 정수 \(381\)에 극히 가깝다. 정확히 \(381\) 단계로
	\[
	\Mpl^{\mathrm{(ladder)}} = m_e \cdot q^{381} \approx 1.219\times10^{19}\ \mathrm{GeV},
	\]
	을 얻으며, 이는 표준값과 단지 \(0.15\%\) 차이난다. 이 놀라운 일치는 숫자 \(2/3\)의 기본적 역할을 추가로 지지한다. \textbf{상태:} 현상학적 관찰.
	
	\section{창발적 시공간 계량 — 추측}
	\label{sec:metric}
	시공간 간격이 사전의 양자 정보 계량으로부터 유도될 수 있다고 추측한다:
	\[
	g_{\mu\nu}(x) \, dx^\mu dx^\nu = 1 - \bigl| \bra{\psi(x)} \psi(x+dx) \rangle \bigr|^2,
	\]
	여기서 \(\ket{\psi}\)는 국소 조정 상태이다. \textbf{상태:} 추측.
	
	\section{조정 장에 대한 유효 스칼라-텐서 작용}
	\label{sec:action}
	\(\phi = \ln(\mathcal{K}_{\mathrm{eff}}/K_0)\)라 하자. 등가 원리를 존중하는 가장 간단한 공변 작용은
	\begin{equation}
		S = \int d^4x\sqrt{-g}\left[ \frac{1}{16\pi G_0} R + \frac12(\nabla\phi)^2 - V(\phi) + \frac{\alpha_2}{2}R\phi^2 \right],
		\label{eq:action}
	\end{equation}
	여기서 \(V(\phi) = V_0 e^{-\lambda\phi}\), \(\lambda = 3/2\)이다. \textbf{상태:} 가정 (이 작용은 첫 원리로부터 유도된 것이 아니라 잘 동기 부여된 유효 기술이다). 상수 \(G_0, \alpha_2, V_0\)는 현상학적이며 관측에 맞추어져야 한다.
	
	\subsection{장 방정식}
	\(g^{\mu\nu}\)와 \(\phi\)에 대한 변분은 다음을 제공한다:
	\begin{align}
		G_{\mu\nu} &= 8\pi G_{\mathrm{eff}}\left(T_{\mu\nu}^{(\phi)} + T_{\mu\nu}^{(\text{물질})}\right) + \frac{\alpha_2}{1-4\pi G_0\alpha_2\phi^2}\left(\nabla_\mu\nabla_\nu\phi^2 - g_{\mu\nu}\Box\phi^2\right), \label{eq:einstein-mod} \\
		G_{\mathrm{eff}} &= \frac{G_0}{1-4\pi G_0\alpha_2\phi^2},\quad
		T_{\mu\nu}^{(\phi)} = \nabla_\mu\phi\nabla_\nu\phi - g_{\mu\nu}\bigl(\tfrac12(\nabla\phi)^2 + V(\phi)\bigr).
	\end{align}
	극한 \(\phi\to0\)에서 일반 상대성 이론이 회복된다. \textbf{상태:} 작용이 주어졌을 때의 정리.
	
	%%%%%%%%%%%%%%%%%%%%%%%%%%%%%%%
	% 제2부: 양자역학과 온도 효과
	%%%%%%%%%%%%%%%%%%%%%%%%%%%%%%%
	\part{양자역학과 온도 효과}
	
	\section{유한 채널 용량으로부터의 양자 불확정성}
	란다우어 경계와 에너지-시간 불확정성으로부터
	\[
	\Delta S \, \Delta K_{\mathrm{eff}} \ge \frac{\hbar}{2}.
	\]
	\textbf{상태:} 그럴듯한 논증; 완전한 유도는 사전의 미시적 모델을 필요로 한다.
	
	\section{중력의 온도 의존성}
	\label{sec:temperature}
	유효 중력 상수는 약한 온도 의존성을 얻는다:
	\[
	\frac{\Delta G}{G} \sim 10^{-16}\,\frac{\Delta T}{T_c},
	\]
	이는 정밀 기기로 테스트 가능하다. \textbf{상태:} 예측.
	
	%%%%%%%%%%%%%%%%%%%%%%%%%%%%%%%
	% 제3부: 우주론과 입자 계층
	%%%%%%%%%%%%%%%%%%%%%%%%%%%%%%%
	\part{우주론과 입자 계층}
	
	\section{우주 상수 (추측)}
	\label{sec:Lambda}
	\(\Omega_\Lambda = 2/3\)는 19차원 전-기하학의 위상에서 발생한다고 추측된다. \textbf{상태:} 추측.
	
	\section{암흑 물질로서의 조정 장 효과}
	\label{sec:DM}
	\eqref{eq:einstein-mod}로부터의 유효 암흑 물질 밀도는 평평한 회전 곡선과 \(\Omega_{\coord} \approx 0.283\)을 제공한다. \textbf{상태:} 유효 작용 가설 하의 정리.
	
	\section{페르미온 질량 계층 (현상학적)}
	\label{sec:masses}
	\[
	m_f = M_{\mathrm{Pl}} \left(\frac{2}{3}\right)^{96 + k_f} \cdot \xi_f,
	\]
	여기서 \(k_f\)와 \(\xi_f\)는 관측에 맞춰진다; \(k_f\)의 반정수 패턴은 위상 기원을 암시한다. \textbf{상태:} 현상학적 법칙.
	
	\section{양자 비국소성 (추측)}
	\label{sec:nonlocality}
	얽힘은 19차원 섬유에서의 기하학적 사전-조정으로 해석된다. \textbf{상태:} 추측.
	
	% ============================================================
	% 새로운 수정 섹션: 소수
	% ============================================================
	\section{소수 조정 사다리 (정리와 경험적 도구)}
	\label{sec:primes}
	
	보편 오차 법칙 \(\varepsilon = \kappa_c\alpha K_{\mathrm{eff}}^{-\beta}\) (\(\beta=2/3\), \(\kappa_c=1/3\))는 모든 조정 시스템에 적용된다. 예상치 못한 귀결은 소수의 분포도 동일한 스케일링을 반영한다는 것이다. YUCT 그림에서 각 정수는 조정 사전의 가능한 상태이며, 소수는 더 간단한 사전 활성화 항목으로 분해될 수 없는 상태이다 — 그것은 최대 조정 완전성을 가진다.
	
	가설 \(K_{\mathrm{eff}}(p_n) \propto p_n\) (더 큰 정수는 비례하여 더 큰 사전을 필요로 함) 하에서 보편 오차 법칙은
	\[
	\varepsilon(p_n) \propto p_n^{-2/3}.
	\]
	이 된다. 이 예측은 고전적 Rosser 근사 \(R_n\) (\(n\)번째 소수에 대한 가장 잘 알려진 기초적 추정)에 대해 테스트된다. \(n = 5\times10^5\)까지의 수치 분석은 국소 스케일링 지수 \(\gamma\)가 \(n\)의 증가와 함께 \(2/3\)로 수렴함을 보여주며 (점근값 \(0.66666\)), 이는 \(\beta\)의 보편성에 대한 독립적 확인을 제공한다.
	
	\subsection{이론적 YUCT 소수 보정 (정리)}
	
	YUCT의 대수 루프 (\(S_{\mathrm{odd}}=1.2\), \(S_{\mathrm{even}}=0.8\))로부터 무매개변수 점근 공식을 얻는다:
	\[
	\boxed{p_n \approx R_n - \frac{S_{\mathrm{even}}}{2}\, n^{1-\beta}\,\ln n},
	\qquad \beta = \frac{2}{3}.
	\]
	이것은 본문의 정리~4이다. 소수 변동의 매끄러운 포락선을 포착한다: 진정한 소수에 대한 상대 오차는 \(n\to\infty\)에서 0으로 간다. 작은 \(n\) (\(n < 1000\))에 대해 이 공식은 리만 \(\zeta\) 함수의 영점으로부터의 진동적 기여를 무시하기 때문에 수치 정확도는 제한적이다.
	
	\subsection{경험적 정제 도구}
	
	유한 범위에서의 실용적 계산을 위해 경험적으로 교정된 보정도 제공한다:
	\[
	p_n \approx R_n + A\, n^{1-\beta}\,(\ln n)^B,
	\qquad A \approx -0.44,\quad B \approx 1.05,
	\]
	이는 \(10^3\le n\le 10^5\)에서의 최소 제곱 피팅으로 얻어진다. 이 도구는 교정 구간에서 최대 오차를 \(\approx 0.006\%\)로 줄인다. 최종 근접 탐색이 정확한 소수 식별을 보장한다. 상수 \(A\)와 \(B\)는 YUCT로부터 유도되지 않으며 순수히 경험적으로 간주되어야 한다; 교정 범위 밖에서의 유효성은 보장되지 않는다.
	
	\subsection{잔여 진동과 \(\zeta(s)\)의 영점}
	
	진정한 소수와 이론적 YUCT 공식의 차이는 리만 \(\zeta\) 함수의 비자명 영점의 기여에 의해 거의 완벽하게 (\(R^2>0.99\)) 기술되는 결정론적 진동으로 구성된다. 부록 PrimeN에서 보고된 이 관찰은 그 영점들이 조정 장의 고유 진동수임을 강력히 시사한다. 탐색 없이 \(p_n\)에 대한 완전한 무매개변수 닫힌 형식 표현은 YUCT 기반 영점 유도를 필요로 하며, 이는 향후 연구로 남겨진다.
	
	\textbf{결과의 상태:}
	\begin{itemize}
		\item 이론적 YUCT 공식: \textbf{정리} (점근적).
		\item 경험적 YUCT 공식: \textbf{경험적 도구}, 정리가 아님.
		\item \(\zeta(s)\) 영점과의 스펙트럼 연관성: \textbf{경험적 관찰}, 이론적 설명은 미해결.
	\end{itemize}
	
	%%%%%%%%%%%%%%%%%%%%%%%%%%%%%%%
	% 제4부: 19차원 모체 이론
	%%%%%%%%%%%%%%%%%%%%%%%%%%%%%%%
	\part{19차원 모체 이론}
	
	\section{19차원 작용의 명시적 형태}
	\label{sec:19Dlagrangian}
	\[
	S_{19} = \frac{1}{2\kappa_{19}} \int d^{19}X \sqrt{-G}\; R(G) \;+\; \int d^{19}X \sqrt{-G}\; \mathcal{L}_{\text{coord}},
	\]
	여기서
	\[
	\mathcal{L}_{\text{coord}} = -\frac{1}{4} \Tr(\Psi_{MN}\Psi^{MN}) + \frac{1}{2} G^{MN}\,\partial_M K_{\mathrm{eff}}\,\partial_N K_{\mathrm{eff}} - V(K_{\mathrm{eff}}) + \dots
	\]
	\textbf{상태:} 공준.
	
	%%%%%%%%%%%%%%%%%%%%%%%%%%%%%%%
	% 제5부: 검증 가능한 예측과 향후 계획
	%%%%%%%%%%%%%%%%%%%%%%%%%%%%%%%
	\part{검증 가능한 예측과 향후 계획}
	
	\section{구체적 실험 테스트}
	\begin{itemize}
		\item \textbf{텐서-스칼라 비율:} \(r \approx 0.07\).
		\item \textbf{서브마이크론 중력:} 10 nm에서 \(\Delta g/g \sim 10^{-12}\).
		\item \textbf{CHSH 부등식:} \(2\sqrt{2}\)를 초과하는 위반 없음.
		\item \textbf{우주 예산:} \(\Omega_\Lambda = 2/3\), \(\Omega_{\coord} = 0.283\), \(\Omega_b = 0.05\).
		\item \textbf{\(G\)의 온도 의존성:} \(\Delta G/G \sim 10^{-16}\).
		\item \textbf{소수 점근 스케일링:} 국소 지수 \(\gamma\)는 \(n\to\infty\)에서 \(2/3\)로 수렴해야 함 (\(n\le 5\times10^5\)에서 이미 검증됨).
	\end{itemize}
	
	\section{완전한 유도를 위한 계획}
	\begin{enumerate}
		\item 내부 다양체 \(X_{15}\)의 구축과 페르미온 질량 매개변수 계산.
		\item 19차원 이론으로부터의 유효 작용 유도.
		\item 사전의 제일원리 통계 역학.
		\item 조정 장의 스펙트럼 특성으로부터 \(\zeta(s)\)의 비자명 영점 유도.
	\end{enumerate}
	
	\section{결론}
	우리는 명시적 비용 함수 분석과 최적 수송 정리를 포함하여 YUCT의 공리로부터 보편 지수 \(\beta = 2/3\)의 자기-일관된 유도를 제시했다. 동일한 지수가 소수의 점근 분포를 지배함이 보여졌으며, 이는 이론을 강화할 뿐만 아니라 수론과 조정 물리학 사이에 새로운 다리를 연다. 이론의 나머지 요소들은 추측 또는 유효 기술로 명확히 표시되어, 추가 연구를 위한 견고한 기초를 제공한다.
	
	\begin{center}
		\textit{우주는 자기-조정하는 사전이다.}
	\end{center}
	
	\begin{thebibliography}{99}
		\bibitem{YakushevLaw2026}
		A.\,V.\,Yakushev, \emph{Yakushev's Law of Coordination: The Fundamental Law of Reality as a Fractal Coordination Network}, Zenodo (2026), DOI:10.5281/zenodo.18444598.
		\bibitem{AppendixL}
		A.\,V.\,Yakushev, ``Appendix L: Fractal Coordination Error Scaling — A Universal Law from DNA to Cosmology,'' in \emph{YUCT}, Zenodo (2026).
		\bibitem{AppendixY}
		A.\,V.\,Yakushev, ``Appendix Y: Mathematical Foundations of YUCT — A Derivation from First Principles,'' in \emph{YUCT}, Zenodo (2026).
		\bibitem{AppendixPrimeN}
		A.\,V.\,Yakushev, ``Appendix PrimeN: YUCT Prime Numbers Coordination Ladder,'' in \emph{YUCT}, Zenodo (2026).
		\bibitem{Rosser1941}
		J.\,B.\,Rosser, ``The \(n\)-th Prime,'' \emph{Amer. Math. Monthly} \textbf{48}, 607--611 (1941).
		\bibitem{West1997}
		West, G. B., Brown, J. H., \& Enquist, B. J. (1997). ``A general model for the origin of allometric scaling laws in biology.'' \emph{Science} \textbf{276}, 122--126.
		\bibitem{Banavar1999}
		Banavar, J. R., Maritan, A., \& Rinaldo, A. (1999). ``Size and form in efficient transportation networks.'' \emph{Nature} \textbf{399}, 130--132.
	\end{thebibliography}
	
\end{document}